% Source: https://artofproblemsolving.com/community/c3h2427111_several_problems_that_i_made
% Stable identity: aops-topic-2427111-q06
There are $\dbinom92=36$ equally likely pairs of cards for each draw, and the two draws are independent.

\textbf{Alice.} The prime sums possible are $3,5,7,11,13,17$:
\begin{itemize}
\item $3$: $\{1,2\}$;\quad $5$: $\{1,4\},\{2,3\}$;\quad $7$: $\{1,6\},\{2,5\},\{3,4\}$;
\item $11$: $\{2,9\},\{3,8\},\{4,7\},\{5,6\}$;\quad $13$: $\{4,9\},\{5,8\},\{6,7\}$;\quad $17$: $\{8,9\}$.
\end{itemize}
That is $14$ pairs, so $P(\text{Alice wins})=\dfrac{14}{36}=\dfrac7{18}$.

\textbf{Bob.} Sums $12$: $\{3,9\},\{4,8\},\{5,7\}$; $13$: $3$ pairs (above); $14$: $\{5,9\},\{6,8\}$; $15$: $\{6,9\},\{7,8\}$; $16$: $\{7,9\}$; $17$: $\{8,9\}$. That is $12$ pairs, so $P(\text{Bob wins})=\dfrac{12}{36}=\dfrac13$.

By independence, the probability that both win is
\[\dfrac7{18}\cdot\dfrac13=\boxed{\dfrac{7}{54}}.\]
